\documentclass[10pt,a4paper]{amsart}
\usepackage[margin=19mm]{geometry}
\usepackage{amsmath,amssymb,mathtools}
\usepackage[scr=rsfs,cal=euler]{mathalfa}
\usepackage{xfrac}
\usepackage[unicode,hidelinks,bookmarks=false]{hyperref}
\newcommand{\ie}{i.e.\ }
\newcommand{\cf}{cf.\ }
\newcommand{\resp}{respectively\ }
\newcommand{\Subsection}{\subsection}
\providecommand{\nobreakdash}{\nobreak\hbox{-}}
\setlength{\emergencystretch}{2em}
\multlinegap0pt
\usepackage{lmodern}
\usepackage{mathtools}
\usepackage{booktabs,array}
\usepackage{flafter}
\hypersetup{pdftitle={Finite deletion-induced saturation for every non-complete graph},pdfauthor={Haochen Liu},pdfsubject={Graph theory; computer-assisted proof}}
\newtheorem{theorem}{Theorem}[section]
\newtheorem{lemma}[theorem]{Lemma}
\newtheorem{proposition}[theorem]{Proposition}
\newtheorem{corollary}[theorem]{Corollary}
\newtheorem{definition}[theorem]{Definition}
\newtheorem{remark}[theorem]{Remark}
\newcommand{\comp}[1]{\overline{#1}}
\newcommand{\Rclass}{\mathcal{R}}
\newcommand{\Dsat}{\mathsf{D}}
\newcommand{\Asat}{\mathsf{A}}
\newcommand{\Del}{\mathrm{del}}
\newcommand{\Add}{\mathrm{add}}
\newcommand{\Sym}{\mathbin{\triangle}}
\newcommand{\id}{\operatorname{id}}
\newcommand{\dom}{\operatorname{dom}}
\newcommand{\wt}{\operatorname{wt}}
\newcommand{\file}[1]{\texttt{\detokenize{#1}}}
\setcounter{tocdepth}{1}
\begin{document}
\title[Finite deletion-induced saturation for every non-complete gr]{Finite deletion-induced saturation for every non-complete graph}
\author{Haochen Liu}
\date{}
\maketitle
\noindent\emph{Source and licence.} Finite deletion-induced saturation for every non-complete graph. Version arXiv:2609.21388v1 (CC BY 4.0). This material is distributed under the Creative Commons Attribution 4.0 International licence, http://creativecommons.org/licenses/by/4.0/, as recorded in the arXiv OAI licence metadata for this exact version and shown on the abstract page https://arxiv.org/abs/2609.21388v1. Full rights notice: licenses/arxiv-2609.21388-CC-BY-4.0.txt.\par\smallskip
\noindent\textit{Editorial scope.} Counted unit: the original \texttt{theorem} environment \texttt{thm:deletion-criterion} at source lines 199--201 together with its complete proof at lines 202--210; the delivered document keeps source lines 119--211 so that every referenced label and every citation resolves inside it. Native editable source is the paper's own concatenated TeX; the AMS wrapper is editorial formatting only. No publisher class, upstream script or unrelated artwork is redistributed.
\begin{theorem}[Finite local lifting, {\cite[Theorem~4.19]{ABO}}]\label{thm:lifting}
For each positive integer $n$, there is a finite extension $B=A_PQ$ such that the seed $A$ embeds as an induced substructure, the specified partial automorphisms extend to automorphisms of $B$, and the natural map
\[
 \pi:T=A_PF\longrightarrow B=A_PQ
\]
is a relational homomorphism. Every weak substructure $U\subseteq B$ of weight at most $n$ admits a homomorphism $\phi:U\to T$ with
\begin{equation}\label{eq:right-inverse}
 \pi\phi=\id_U.
\end{equation}
Every relation tuple of $B$ lies in a $Q$-translate of a seed relation tuple.
\end{theorem}

Here a weak substructure may retain only some relation tuples supported on its vertex set. We apply the theorem to the full induced relational structure on that set.

\begin{lemma}[Induced lifting]\label{lem:induced-lifting}
If $U$ is the full induced substructure of $B$ on a vertex set of weight at most $n$, the map $\phi$ in Theorem~\ref{thm:lifting} is an induced embedding.
\end{lemma}
\begin{proof}
Equation~\eqref{eq:right-inverse} makes $\phi$ injective, and its homomorphism property preserves every relation tuple of $U$. If a tuple of elements of $\phi(U)$ belonged to a relation in $T$ while its preimage tuple did not belong to that relation in $U$, the homomorphism $\pi$ would send it to the missing tuple in $B$. This contradicts the assumption that $U$ contains all relations induced by $B$ on its vertices.
\end{proof}

For our applications, at most two binary relations are present and at most $h$ vertices must be lifted. If $r$ vertices occur in no relation tuple, then
\[
 \wt(U)\le r+2(h-r)^2\le 2h^2\qquad(h\ge1).
\]
This bound allows both loops and overlaps between the relations before they have been ruled out. We may therefore use
\begin{equation}\label{eq:weight-bound}
 n=2h^2+4.
\end{equation}
Once the free extension is known to be loopless and to have disjoint $E$ and $N$, lifting substructures on one or two vertices shows the same properties for $B$. More generally, an induced graph on at most $h$ vertices in the underlying $E$-graph of $B$ lifts inducedly to the underlying $E$-graph of $T$, by retaining the full induced structure in all relation symbols before applying the lemma.

\section{Two-point amalgamation and finite protection}\label{sec:criteria}

\subsection{The tree of seed copies}

Let $A$ be a finite graph and let $e_0$ be a specified adjacent pair or a specified nonadjacent pair. For every other pair $f$ of the same type, choose a bijection $p_f:e_0\to f$. These are partial automorphisms. When nonedges are protected, every seed pair has exactly one of the two relations $E,N$, and $p_f$ preserves both relations.

The left Cayley graph of the free group on the symbols $p_f$ is a tree, with edges joining $g$ to $p_fg$. In the free extension, the corresponding seed copies $A_g$ and $A_{p_fg}$ are identified along the two endpoints prescribed by $p_f$. The following properties justify the separator arguments used below.

\begin{lemma}\label{lem:tree-structure}
The seed copies in this two-point free extension satisfy the following properties.
\begin{enumerate}
\item Every $A_g$ embeds inducedly; distinct vertices within a seed copy are not identified.
\item The set of bags containing a fixed quotient vertex is a connected subtree of the index tree.
\item Cutting an edge of the index tree separates its bags into two sides whose vertex sets intersect only in the two-vertex seam. No graph edge joins vertices exclusive to opposite sides.
\item In the two-relation construction, $E$ and $N$ are irreflexive and disjoint throughout the free extension.
\end{enumerate}
\end{lemma}
\begin{proof}
An identification chain projects to a walk in the index tree. A chain starting and ending in the same bag projects to a closed walk. Successively cancelling backtracking pairs in this walk cancels mutually inverse partial bijections, so its initial and final coordinates agree. Thus no bag acquires identifications between distinct vertices.

If a quotient vertex lies in two bags, an identification chain joins those occurrences. Cancelling backtracking shows that it occurs in every bag along the unique tree path between them. This proves the connectedness assertion and also shows that a vertex appearing on both sides of a cut must belong to its seam.

Every graph edge comes from a single bag. Consequently an edge cannot join vertices exclusive to opposite sides. To see that a bag receives no additional internal relation, suppose a different bag supplies a relation between two of its vertices. Both vertices occur along the entire path joining the two bags. They therefore cross each seam together. The two-point identifications preserve the relation on that seam, so the relation was present in the original bag as well. The same argument for each of $E$ and $N$ proves inducedness and excludes conflicts; the absence of internal identifications excludes loops.
\end{proof}

\begin{lemma}[Bag exclusion]\label{lem:bag-exclusion}
Suppose $C$ is connected, has no cut vertex, and has no two-vertex cut of the seam type. If $A$ is $C$-free, then its two-point free extension is $C$-free.
\end{lemma}
\begin{proof}
Suppose an induced copy of $C$ exists in the extension. Choose a finite connected subtree with the fewest bags whose union covers its vertices. A single bag contradicts the $C$-freeness of $A$. Otherwise, cut the tree edge incident with a leaf bag. Minimality implies that the copy has a vertex exclusive to the leaf side and a vertex exclusive to the opposite side. Its intersection with the seam is therefore a separating set. By Lemma~\ref{lem:tree-structure}, this set has size zero, one, or two; in the last case its type agrees with the seam. These alternatives contradict, respectively, connectedness, the absence of a cut vertex, and the assumed exclusion of the relevant two-vertex cut. The separation is in the full index tree, so bags outside the chosen finite subtree cannot provide an edge between the two sides.
\end{proof}

\subsection{Completing protected nonedges}

\begin{lemma}[Protected completion]\label{lem:completion}
Let $E$ and $N$ be disjoint sets of unordered pairs on a finite set $V$, with $N\ne\varnothing$. Suppose $(V,E)$ is $H$-free. Assume that each $f\in N$ has a witness on a vertex set $S_f$ for which every pair is specified in $E$ or $N$, and adding $f$ to the graph on $S_f$ gives an induced copy of $H$. Then $\Asat(H)$ holds.
\end{lemma}
\begin{proof}
Choose a graph $J$ with the maximum number of edges among all $H$-free graphs satisfying
\[
 E\subseteq E(J)\subseteq\binom V2\setminus N.
\]
The set of choices is finite and nonempty. If a nonedge $f$ of $J$ lies in $N$, all pairs of the specified witness $S_f$ retain their original status, so $J+f$ contains an induced $H$. If $f\notin N$, then $J+f$ cannot be $H$-free, by maximality. Finally, $N\ne\varnothing$ ensures that $J$ is not complete.
\end{proof}

The argument does not require induced-$H$-freeness to be monotone under edge addition. An implementation which adds allowable edges must continue until none remain allowable; a single scan need not suffice.

\subsection{The two criteria}


\begin{theorem}[Deletion criterion]\label{thm:deletion-criterion}
Suppose $H$ contains a connected induced subgraph $C$ having neither a cut vertex nor an adjacent two-vertex cut. If $H$ has a nonedge $e_0$ such that $H+e_0$ is $C$-free, then $\Dsat(H)$ holds.
\end{theorem}

\begin{proof}
Put $A=H+e_0$ and map the marked edge $e_0$ to each other edge of $A$ by a two-point partial automorphism. The resulting free extension $T$ is $C$-free by Lemma~\ref{lem:bag-exclusion}. Apply Theorem~\ref{thm:lifting} with the bound~\eqref{eq:weight-bound}, where $h=|V(H)|$. An induced $C$ in the finite extension $B$ would lift inducedly to $T$, by Lemma~\ref{lem:induced-lifting}. Thus $B$ is $C$-free and hence $H$-free.

Every edge $e$ of $B$ is a translate of a seed edge $f$. The extended partial automorphism taking $e_0$ to $f$, followed by that translation, is an automorphism $\alpha$ of $B$ taking $e_0$ to $e$. For $f=e_0$, use the identity in the first step. Since $A$ embeds inducedly,
\[
 (B-e)[\alpha(V(A))]\cong A-e_0=H.
\]
This proves the required property for every edge of $B$. The seed edge $e_0$ ensures $E(B)\ne\varnothing$. If $A$ has only one edge, it already is the required host, so the construction with no generators causes no exception.
\end{proof}

\begin{thebibliography}{99}
\bibitem[ABO]{ABO} K.~Auinger, J.~Bitterlich and M.~Otto, \emph{Finite approximation of free groups II: The Theorems of Ash, Herwig--Lascar and Ribes--Zalesskii---revisited and strengthened}, preprint, arXiv:2507.11685v5 (2026). \url{https://arxiv.org/abs/2507.11685v5}.
\end{thebibliography}
\end{document}
